
\documentclass[a4paper,11pt]{article}
\usepackage{amsmath,amsthm,amsfonts,amssymb,amscd,amstext,vmargin,graphics,graphicx,tabularx,multicol} \usepackage[french]{babel}
\usepackage[utf8]{inputenc}  
\usepackage[T1]{fontenc} 
\usepackage[T1]{fontenc}
\usepackage{amsmath,amssymb}
\usepackage{pstricks-add,tikz,tkz-tab,variations}
\usepackage[autolanguage,np]{numprint} 
\usepackage{color}
\usepackage{ulem}

\setmarginsrb{1.5cm}{0.5cm}{1cm}{0.5cm}{0cm}{0cm}{0cm}{0cm} %Gauche, haut, droite, haut
\newcounter{numexo}
\newcommand{\exo}[1]{\stepcounter{numexo}\noindent{\bf Exercice~\thenumexo} : \marginpar{\hfill /#1}}
\reversemarginpar


\newcounter{enumtabi}
\newcounter{enumtaba}
\newcommand{\q}{\stepcounter{enumtabi} \theenumtabi.  }
\newcommand{\qa}{\stepcounter{enumtaba} (\alph{enumtaba}) }
\newcommand{\initq}{\setcounter{enumtabi}{0}}
\newcommand{\initqa}{\setcounter{enumtaba}{0}}

\newcommand{\be}{\begin{enumerate}}
\newcommand{\ee}{\end{enumerate}}
\newcommand{\bi}{\begin{itemize}}
\newcommand{\ei}{\end{itemize}}
\newcommand{\bp}{\begin{pspicture*}}
\newcommand{\ep}{\end{pspicture*}}
\newcommand{\bt}{\begin{tabular}}
\newcommand{\et}{\end{tabular}}
\renewcommand{\tabularxcolumn}[1]{>{\centering}m{#1}} %(colonne m{} centrée, au lieu de p par défault) 
\newcommand{\tnl}{\tabularnewline}

\newcommand{\trait}{\noindent \rule{\linewidth}{0.2mm}}
\newcommand{\hs}[1]{\hspace{#1}}
\newcommand{\vs}[1]{\vspace{#1}}

\newcommand{\N}{\mathbb{N}}
\newcommand{\Z}{\mathbb{Z}}
\newcommand{\R}{\mathbb{R}}
\newcommand{\C}{\mathbb{C}}
\newcommand{\Dcal}{\mathcal{D}}
\newcommand{\Ccal}{\mathcal{C}}
\newcommand{\mc}{\mathcal}

\newcommand{\vect}[1]{\overrightarrow{#1}}
\newcommand{\ds}{\displaystyle}
\newcommand{\eq}{\quad \Leftrightarrow \quad}
\newcommand{\vecti}{\vec{\imath}}
\newcommand{\vectj}{\vec{\jmath}}
\newcommand{\Oij}{(O;\vec{\imath}, \vec{\jmath})}
\newcommand{\OIJ}{(O;I,J)}

\newcommand{\bmul}[1]{\begin{multicols}{#1}}
\newcommand{\emul}{\end{multicols}}


\newcommand{\reponse}[1][1]{%
\multido{}{#1}{\makebox[\linewidth]{\rule[0pt]{0pt}{20pt}\dotfill}
}}

\newcommand{\titre}[5] 
% #1: titre #2: haut gauche #3: bas gauche #4: haut droite #5: bas droite
{
\noindent #2 \hfill #4 \\
#3 \hfill #5

\vspace{-1.6cm}

\begin{center}\rule{6cm}{0.5mm}\end{center}
\vspace{0.2cm}
\begin{center}{\large{\textbf{#1}}}\end{center}
\begin{center}\rule{6cm}{0.5mm}\end{center}
}



\begin{document}
\pagestyle{empty}
\titre{Devoir maison}{Nom :}{Prénom :}{3 ème}{Date :}

\vspace*{1cm}


\exo

Le graphique suivant donne la température  relevée au village  olympique de Pyeongchang le 31  janvier 2018. 

\q Quelle température faisait-il à 8h ? \textbf{-4 \degre C} \\

\q A quelle(s) heure(s) la température était-elle de -2\degre C ? \textbf{11h et 18h} \\

\q Quelle est la température minimale de la journée et à quelle heure a-t-elle été relevée ? \textbf{-12 \degre C vers 3h} \\

\q Pendant combien de temps la température est-elle restée positive au cours de cette journée ? \textbf{4h30} \\

\q On associe la courbe ci-dessus à une fonction $f$ .  \\

   
  \qa Quelle est l'image de 9 par cette fonction ? \textbf{$h(9)=-4$} \\
  
  \qa Quelle valeur a pour antécédent 14 ? \textbf{$h(3) = 14$}\\
  
  \qa Quelle(s) valeur(s) a/ont pour image -8 ? \textbf{$h(5)=-8$}\\
  
   \qa Quel(s) est (sont) le (les) antécédent(s) de 2 par cette fonction ? \textbf{$h(13)=2 $ et $h(15)=2$ }\\
   
   \vspace*{1cm}


\exo

On donne l'expression $D=(2x+3)(2x-1)-(3x-4)(2x+3)$.\\

\noindent \initq \q $D=(2x+3)(2x-1)-(3x-4)(2x+3)$ \\
$D= 4x^{2}-2x+6x-3 - (6x^{2} + 9x -8x-12)$\\
$D= 4x^{2}-2x+6x-3 - 6x^{2} - 9x +8x+12$\\
\fbox{$D= -2x^{2} +3x +9$}\\


\q $D=(2x+3)(2x-1)-(3x-4)(2x+3)$\\
$D=(2x+3)[(2x-1)-(3x-4)]$ \\
$D=(2x+3)(2x-1-3x+4)$\\
\fbox{$D=(2x+3)(-x+3)$}\\



\q Pour x = 1, je choisi l'expression factorisée de D de la question 2.\\
$D=(2 \times 1 +3)(-1+3)$\\
$D=5\times 2$\\
\fbox{$D=10$}\\

\q Pour $x=-\dfrac{3}{2}$, je choisi l'expression développée de D de la question 1.\\
$D= -2 \times \left(-\dfrac{3}{2}\right)^{2} +3\times\left(-\dfrac{3}{2}\right) +9$\\
$D= -2 \times \left(\dfrac{9}{4}\right) +3\times\left(-\dfrac{3}{2}\right) +9$\\
$D= -\dfrac{9}{2} -\dfrac{9}{2} +9$\\
$D= -\dfrac{18}{2} +9$\\
\fbox{$D= 0$}\\

\vspace*{1cm}

\newpage

\exo

La figure suivante est donnée à titre indicatif pour préciser la position des points A, B, C, D et E. Les longueurs représentées ne le sont pas en vraie grandeur.\\
On donne CE = 5 cm, CD = 12 cm, CA = 18 cm, CB = 7,5 cm, AB = 19,5 cm.

\begin{center}
\includegraphics[scale=1]{dmexo3.eps} 
\end{center}

\noindent \initq \q Les points E, C et B sont alignés dans le même ordre que les points D, C et A.

\bmul{2}
D'une part, $\dfrac{EC}{CB}=\dfrac{5}{7.5}\approx 0,67$


\columnbreak
D'autre part, $\dfrac{DC}{CA}=\dfrac{12}{18} \approx 0,67$

\emul

Ainsi $\dfrac{EC}{CB} = \dfrac{DC}{CA}$, donc d'après la réciproque du théorème de Thalès les droites (ED) et AB) sont parallèles.\\


\q Dans les triangles EDC et ABC :\\

- Les points E, C et B sont alignés dans le même ordre que les points D, C et A.\\
- Les droites (ED) et AB) sont parallèles\\

D'après le théorème de Thalès, on a :  $\dfrac{EC}{CB}=\dfrac{DC}{CA} = \dfrac{ED}{AB}$\\
\begin{center}
$\dfrac{5}{7.5}= \dfrac{12}{18} = \dfrac{ED}{19.5}$\\
\end{center}

$ \dfrac{12}{18} = \dfrac{ED}{19.5}$ donc $ED = \dfrac{12 \times 19.5}{18}=13 cm$\\





\q  Dans le triangle DCE,\\

\bmul{2}
D'une part, $ED^{2}=13^{2} = 169$


\columnbreak
D'autre part, $EC^{2} + CD^{2} = 5^{2} + 12^{2} = 25 + 144 $\\
\hspace*{1cm} $EC^{2} + CD^{2} =169$

\emul

Ainsi, $EC^{2} + CD^{2} =ED^{2}$ donc d'après la réciproque du théorème de Pythagore le triangle CDE est rectangle en C.\\




\q Dans le triangle DCE rectangle en C, je cherche l'angle $\widehat{DEC}$, je connais l'hypoténuse et le côté adjacent.\\

\hspace*{2cm} \noindent $cos \widehat{DEC} =\dfrac{EC}{ED} $\\
\hspace*{2.5cm} $cos \widehat{DEC} =\dfrac{5}{13} $\\

A l'aide de la calculatrice, on trouve $\widehat{DEC} \approx 67  $ \degre\\

\end{document}
